\chapter{The Spectral Theorem and Rayleigh Quotients}
\label{ch:spectral-theorem-rayleigh}

\section{Why Symmetric Matrices Are Special}

General square matrices may have too few eigenvectors, complex eigenvalues, or
badly conditioned eigenvector bases. Symmetric real matrices behave much
better.

\begin{theorem}[Spectral theorem]
Let $A\in\R^{n\times n}$ be symmetric. Then there is an orthonormal basis
$q_1,\ldots,q_n$ of $\R^n$ and real numbers
$\lambda_1,\ldots,\lambda_n$ such that
\[
  Aq_i=\lambda_i q_i.
\]
Equivalently, if $Q=[q_1\ \cdots\ q_n]$, then
\[
  Q^TQ=I,
  \qquad
  A=Q\Lambda Q^T,
  \qquad
  \Lambda=\diag(\lambda_1,\ldots,\lambda_n).
\]
\end{theorem}

This is diagonalization with an orthonormal change of basis. The inverse of
$Q$ is just $Q^T$, so the geometry is preserved.

\section{Orthogonal Eigenvectors}

One piece of the spectral theorem is easy to prove and useful on its own.

\begin{proposition}
Let $A$ be symmetric. If $Av=\lambda v$ and $Aw=\mu w$ with
$\lambda\neq \mu$, then $v\perp w$.
\end{proposition}

\begin{proof}
Using symmetry,
\[
  \lambda\inner{v}{w}
  =
  \inner{\lambda v}{w}
  =
  \inner{Av}{w}
  =
  \inner{v}{Aw}
  =
  \inner{v}{\mu w}
  =
  \mu\inner{v}{w}.
\]
Thus $(\lambda-\mu)\inner{v}{w}=0$. Since $\lambda\neq\mu$, we get
$\inner{v}{w}=0$.
\end{proof}

\begin{computebox}
This is why numerical linear algebra uses \texttt{np.linalg.eigh} for
symmetric matrices rather than the more general \texttt{np.linalg.eig}. The
symmetric problem has real eigenvalues and an orthonormal eigenbasis, and
algorithms exploit that structure.
\end{computebox}

\section{A Useful Proof Idea}

The full spectral theorem is deeper than the proposition above, but the main
mechanism is approachable:

\begin{enumerate}
  \item A symmetric matrix has at least one real eigenvector.
  \item The orthogonal complement of an invariant subspace is invariant.
  \item Peel off one unit eigenvector at a time.
\end{enumerate}

The key invariance statement is short.

\begin{lemma}
Let $A$ be symmetric, and let $U\subseteq\R^n$ be a subspace such that
$A(U)\subseteq U$. Then
\[
  A(U^\perp)\subseteq U^\perp.
\]
\end{lemma}

\begin{proof}
Take $w\in U^\perp$. To show $Aw\in U^\perp$, we must show
$\inner{Aw}{u}=0$ for every $u\in U$. Since $A$ is symmetric,
\[
  \inner{Aw}{u}=\inner{w}{Au}.
\]
Because $U$ is invariant, $Au\in U$. Since $w\in U^\perp$, the last inner
product is zero.
\end{proof}

Once one eigenvector $q_1$ is found, the line $\Span(q_1)$ is invariant, so
its orthogonal complement is invariant. Restrict $A$ to that complement and
repeat. This creates an orthonormal eigenbasis.

\section{Rayleigh Quotients}

\begin{definition}[Rayleigh quotient]
For a symmetric matrix $A$ and nonzero vector $x$, the \emph{Rayleigh quotient}
is
\[
  R_A(x)=\frac{x^TAx}{x^Tx}.
\]
\end{definition}

The quotient is scale-invariant:
\[
  R_A(cx)=R_A(x)
  \qquad(c\neq 0).
\]
So it is enough to study it on unit vectors.

\begin{theorem}[Rayleigh quotient bounds]
Let $A$ be symmetric with eigenvalues
\[
  \lambda_1\leq \lambda_2\leq \cdots \leq \lambda_n.
\]
Then for every nonzero $x$,
\[
  \lambda_1\leq R_A(x)\leq \lambda_n.
\]
Moreover,
\[
  \min_{x\neq 0} R_A(x)=\lambda_1,
  \qquad
  \max_{x\neq 0} R_A(x)=\lambda_n.
\]
\end{theorem}

\begin{proof}
Use the spectral theorem. Write $x$ in the orthonormal eigenbasis:
\[
  x=c_1q_1+\cdots+c_nq_n.
\]
Then
\[
  x^Tx=c_1^2+\cdots+c_n^2
\]
and
\[
  x^TAx=\lambda_1c_1^2+\cdots+\lambda_nc_n^2.
\]
Therefore
\[
  R_A(x)=
  \frac{\lambda_1c_1^2+\cdots+\lambda_nc_n^2}
       {c_1^2+\cdots+c_n^2}.
\]
This is a weighted average of the eigenvalues, with nonnegative weights
$c_i^2$. A weighted average lies between the smallest and largest values.
Equality is achieved by taking $x=q_1$ or $x=q_n$.
\end{proof}

\begin{aimlbox}
For a covariance matrix $C$, the Rayleigh quotient
$w^TCw/(w^Tw)$ is the variance of centered data in direction $w$, normalized
by the length of $w$. The top eigenvector is therefore the direction of
maximum variance. This is the linear algebra behind the first principal
component.
\end{aimlbox}

\section{PSD Matrices In Spectral Form}

\begin{proposition}
If $A$ is symmetric, then $A$ is PSD if and only if all eigenvalues of $A$ are
nonnegative.
\end{proposition}

\begin{proof}
If $A$ is PSD and $Av=\lambda v$ with $v\neq 0$, then
\[
  0\leq v^TAv=v^T(\lambda v)=\lambda\norm{v}^2,
\]
so $\lambda\geq 0$.

Conversely, if all eigenvalues are nonnegative, write
$x=c_1q_1+\cdots+c_nq_n$ in an orthonormal eigenbasis. Then
\[
  x^TAx=\lambda_1c_1^2+\cdots+\lambda_nc_n^2\geq 0.
\]
Thus $A$ is PSD.
\end{proof}

\section{Power Iteration}

The Rayleigh quotient explains why repeated multiplication can find dominant
eigenvectors. If
\[
  x_0=c_1q_1+\cdots+c_nq_n
\]
and $|\lambda_n|$ is strictly larger than the others, then
\[
  A^t x_0
  =
  c_1\lambda_1^tq_1+\cdots+c_n\lambda_n^tq_n
\]
eventually points mostly in the $q_n$ direction, provided $c_n\neq 0$.
Normalizing at each step prevents overflow.

\begin{lstlisting}
import numpy as np

rng = np.random.default_rng(0)
M = rng.standard_normal((4, 4))
A = M + M.T

x = rng.standard_normal(4)
x = x / np.linalg.norm(x)

for _ in range(20):
    x = A @ x
    x = x / np.linalg.norm(x)

rayleigh = x @ A @ x
vals, vecs = np.linalg.eigh(A)

print("Rayleigh quotient:", rayleigh)
print("eigenvalues:", vals)
print("largest abs eigenvalue:", vals[np.argmax(np.abs(vals))])
\end{lstlisting}

Power iteration is a first algorithmic glimpse. Practical eigensolvers are more
careful, but they still exploit the same spectral structure.

\begin{labbox}
Lab 9 experiments with Rayleigh quotients, gradient-style eigenvalue searches,
and deflation. Keep the theorem in view: the quotient is not just a heuristic;
for symmetric matrices its extrema are eigenvalues.
\end{labbox}

\section*{Chapter Summary}
\addcontentsline{toc}{section}{Chapter Summary}

\begin{itemize}
  \item Real symmetric matrices have an orthonormal eigenbasis.
  \item Eigenvectors with distinct eigenvalues of a symmetric matrix are
        orthogonal.
  \item The Rayleigh quotient is a weighted average of eigenvalues.
  \item The minimum and maximum Rayleigh quotients are the smallest and largest
        eigenvalues.
  \item A symmetric matrix is PSD exactly when all its eigenvalues are
        nonnegative.
\end{itemize}

\section*{Exercises}
\addcontentsline{toc}{section}{Exercises}

\subsection*{A. Theory}

\begin{exercise}
\exerciseid{13.A1}
State the spectral theorem for real symmetric matrices and explain why
orthogonality is part of the conclusion.
\end{exercise}

\begin{exercise}
\exerciseid{13.A2}
Let $A$ be symmetric. Prove that eigenvectors with distinct eigenvalues are
orthogonal.
\end{exercise}

\subsection*{B. By Hand}

\begin{exercise}
\exerciseid{13.B1}
Compute the Rayleigh quotient of
\[
  A=\begin{bmatrix}2&0\\0&5\end{bmatrix}
\]
at $x=(1,1)^T$ and at $x=(0,1)^T$.
\end{exercise}

\begin{exercise}
\exerciseid{13.B2}
For the same matrix $A$, find the minimum and maximum possible values of the
Rayleigh quotient.
\end{exercise}

\subsection*{C. Python / NumPy}

\begin{exercise}
\exerciseid{13.C1}
Sample random unit vectors and evaluate the Rayleigh quotient for a symmetric
$2\times 2$ matrix. Compare the observed minimum and maximum with
\texttt{np.linalg.eigvalsh}.
\end{exercise}

\begin{exercise}
\exerciseid{13.C2}
Implement 20 steps of power iteration for a symmetric matrix. Compare the final
Rayleigh quotient with the eigenvalue of largest absolute value.
\end{exercise}

\subsection*{D. Challenge}

\begin{exercise}
\exerciseid{13.D1}
Explain why the maximum Rayleigh quotient of a symmetric matrix should be an
eigenvalue. Use the orthonormal eigenbasis picture.
\end{exercise}

